\section{The unrestricted converse and Goldfeld's densities}
\label{sec:conclusions}

The two residual constructions now give the pointwise theorem. The
density statement then requires only Smith's Selmer distribution and
an elementary change of counting convention.

\begin{proof}[Proof of Theorem~\ref{thm:main-converse}]
By Lemma~\ref{pr:residual-cases}, either $E(\Q)[2]\ne0$ or $E[2]$
is irreducible. Theorem~\ref{thm:rational-converse} proves
$a(E)=c_2(E)$ in the first case, and
Theorem~\ref{thm:irreducible-converse} proves it in the second.
Both use the local Kummer conditions defining the corank in the
statement. The isogeny choices in the rational-torsion proof preserve
that corank and the $L$-function by Lemma~\ref{pr:invariance}.
Thus every curve in the theorem has analytic rank zero or one.
Lemma~\ref{pr:ranks} supplies the Mordell--Weil rank equality and
finiteness of the whole Tate--Shafarevich group.
\end{proof}

\begin{proof}[Proof of Corollary~\ref{cor:finite-descent}]
Lemma~\ref{pr:finite-descent} gives $c_2(E)=d_2(E)$ and shows that
$\Sha(E/\Q)[2^\infty]$ is divisible. Theorem~\ref{thm:main-converse}
gives the rank equalities and whole-Sha finiteness. A finite divisible
$2$-group is zero, so its $2$-primary part vanishes.
\end{proof}

\begin{lemma}[From integer to squarefree parameters]\label{lem:squarefree-density}
Let $\mathcal P$ be a property of nonzero signed integers unchanged by
multiplication by an integer square. Suppose
\[
 A(X)=\#\{n\in\Z:0<|n|\le X,\ \mathcal P(n)\}
      =2\delta X+o(X).
\]
Then
\[
 S(X)=\#\{d\in\mathcal D(X):\mathcal P(d)\}
      =\frac{2\delta}{\zeta(2)}X+o(X).
\]
In particular the property has density $\delta$ among signed
squarefree parameters ordered by absolute value.
\end{lemma}

\begin{proof}
Every nonzero integer has a unique expression $n=m^2d$ with $m>0$
and $d$ signed squarefree. Thus
\[
 A(X)=\sum_{m\le\sqrt X}S(X/m^2),\qquad
 S(X)=\sum_{m\le\sqrt X}\mu(m)A(X/m^2).
\]
The second equality follows by substituting the first and using
$\sum_{m\mid r}\mu(m)=\mathbf1_{r=1}$.
For a fixed integer $M$, the terms with $m\le M$, divided by $X$,
tend to $2\delta\sum_{m\le M}\mu(m)/m^2$. The remaining absolute
contribution is at most
\[
 2X\sum_{m>M}m^{-2},
\]
because $A(Y)\le2\lfloor Y\rfloor\le2Y$.
First let $X$ and then $M$ tend to infinity. The absolutely convergent
series $\sum\mu(m)/m^2=1/\zeta(2)$ proves the formula.
Applying it to the property always true gives
$\#\mathcal D(X)\sim2X/\zeta(2)$, proving the density assertion.
\end{proof}

\begin{proof}[Proof of Theorem~\ref{thm:main-density}]
Smith's theorem states that, for every fixed elliptic curve $E/\Q$
and $j=0,1$ \citep[Theorem~1.1]{Smith2025},
\[
 \#\{n\in\Z:0<|n|\le X,\ c_2(E^{(n)})=j\}=X+o(X)
 .
\]
Quadratic twisting depends only on the squareclass, so
Lemma~\ref{lem:squarefree-density} gives density $1/2$ for each
corank among the signed squarefree parameters.

For $j=0,1$, Theorem~\ref{thm:main-converse} gives
$c_2(E^{(d)})=j\Rightarrow a(E^{(d)})=j$, while
Lemma~\ref{pr:ranks} gives the converse implication.
The analytic-rank counts therefore equal the corresponding corank
counts exactly. The two densities sum to one. Their complement,
the twists of analytic rank at least two, has density zero, and
Lemma~\ref{pr:ranks} gives the final rank and finiteness assertions
on the density-one union.
\end{proof}
